\begin{lemma}[Suffix of a Tashkinov augmentation sequence]
Let (S,B,C) be tree-sequences for the same finite loopless multigraph
with its individual edge type and the same distinguished first edge, and
let \(\varphi\) be one fixed partial edge coloring.  Suppose that \(C\) is
obtained from \(S\) by a
\(\varphi\)-\noderef{TashkinovAugmentationSequence}, that \(S\) is an exact
ordered \noderef{TreeSequenceInitialSegment} of \(B\), and that \(B\) is an
exact ordered \noderef{TreeSequenceInitialSegment} of \(C\).  Then \(C\) is
obtained from \(B\) by a
\(\varphi\)-\noderef{TashkinovAugmentationSequence}.
\end{lemma}
\begin{proof}
We first prove two elementary consequences of exact ordered initial
segments.  If \(X\preccurlyeq Y\) and \(|X|=|Y|\), then \(X=Y\): the
vertex and edge equalities in
\noderef{TreeSequenceInitialSegment} identify the entries at every index,
while proof irrelevance identifies the attachment and injectivity proofs
in the two tree-sequence structures.  Also, if
\(X\preccurlyeq Z\), \(Y\preccurlyeq Z\), and \(|X|\leq |Y|\), then
\(X\preccurlyeq Y\).  Indeed, for each vertex or edge index of \(X\), the
corresponding entries of \(X\) and \(Y\) are both equal to the entry of
\(Z\) at that index; the assumed length inequality supplies the required
index casts.  This verifies all three clauses of
\noderef{TreeSequenceInitialSegment}.

Unfold \noderef{TashkinovAugmentationSequence}.  The given passage from
\(S\) to \(C\) is a finite reflexive-transitive chain of one-step
\(\varphi\)-\noderef{TashkinovAugmentation}s.  We induct on that chain,
allowing \(B\) to vary subject to \(S\preccurlyeq B\preccurlyeq C\), and
prove that its suffix from \(B\) to \(C\) is again such a chain.

If the chain is empty, then \(C=S\).  The two initial-segment length
inequalities give
\[
                  |B|\leq |C|=|S|\leq |B|.
\]
Thus \(|B|=|S|\).  Since \(B\preccurlyeq C=S\), the first elementary
consequence gives \(B=S=C\), and the empty reflexive chain from \(B\) to
\(C\) proves the conclusion.

For the induction step, write the final one-step augmentation as
\(A\longrightarrow C\), preceded by a chain from \(S\) to \(A\).  By
\noderef{TashkinovAugmentation},
\[
                       A\preccurlyeq C,
             \qquad |C|=|A|+1.                                \tag{1}
\]
Suppose first that \(|B|\leq |A|\).  Both \(B\) and \(A\) are exact
ordered initial segments of \(C\), so the second elementary consequence
gives \(B\preccurlyeq A\).  The original relation
\(S\preccurlyeq B\) is unchanged.  The induction hypothesis therefore
gives a \(\varphi\)-augmentation sequence from \(B\) to \(A\).  Appending
the actual last augmentation \(A\to C\), including its same exact new
edge identity, gives the required sequence from \(B\) to \(C\).

It remains that \(|B|>|A|\).  Since \(B\preccurlyeq C\), equation (1)
also gives \(|B|\leq |C|=|A|+1\), so necessarily \(|B|=|C|\).  The first
elementary consequence applied to \(B\preccurlyeq C\) yields \(B=C\).
The empty reflexive chain from \(B=C\) to \(C\) proves the desired suffix.
These two cases complete the induction.
\end{proof}
